\subsection{\#96 --- .osi loader skips all custom-material validation}
\textbf{Link:} \url{https://github.com/Nidhikhare12/OsdagBridge/issues/96}

\textbf{Statement.} The \texttt{.osi} load path validates geometry and additional inputs, but no \texttt{material.*} key is ever validated, so a hand-edited file can carry e.g.\ \texttt{Fy\,=\,99999\,MPa} and load with ``Loaded OSI Successfully!''.

\textbf{Evidence (executed 2026-11-15, twice).}
\begin{verbatim}
validate_osi_inputs({'material.girder.fy': 99999.0,
                     'material.girder.density': 5.0,
                     'material.deck.fck': -3.0})
  -> is_valid = True, errors = []

BridgeInputValidator.validate_material_inputs (the check the
loader never calls) rejects all three:
  fy 99999  -> corrected to 1200.0
  density 5 -> corrected to 50.0
  fck -3    -> corrected to 10.0
\end{verbatim}

\textbf{Root cause.} \texttt{osi\_validator.py}: \texttt{\_BASIC\_KEYS} (L82--86) holds only span/carriageway/skew; \texttt{\_ADDITIONAL\_KEYS} (L89--111) holds many numeric limits but no \texttt{material.*} keys; \texttt{validate()} (L142--157) iterates only those lists, so \texttt{validate\_material\_inputs} --- the check the live Material Properties dialog enforces --- is unreachable from the load path (\texttt{template\_page.py} \texttt{loadOSI\_inputs} L847).

\textbf{Suggested fix (diff posted in the issue).} add the material keys to the OSI validation loop and route them through \texttt{validate\_material\_inputs} in \texttt{validate()}.

\subsection{\#97 --- input lock can crash when Additional Inputs dialog was never opened (T3, static)}
\textbf{Link:} \url{https://github.com/Nidhikhare12/OsdagBridge/issues/97}

\textbf{Statement.} Lock/unlock propagates state to the Additional Inputs dialog, but \texttt{InputDock.toggle\_lock()} dereferences the dialog without checking whether it has been instantiated; because the dialog is created lazily, toggling lock before opening Additional Inputs can raise \texttt{AttributeError} on \texttt{None}. (Reported as a T3 static finding from code reading; not executed as a runtime repro in this session.)

\textbf{Evidence (read-only, file:line).} \texttt{input\_dock.py:856--859} calls \texttt{self.parent.\_additional\_inputs\_dialog.lock(...)} unconditionally, while \texttt{template\_page.py:124} initializes \texttt{\_additional\_inputs\_dialog = None} and creates it on demand (\texttt{\_get\_additional\_inputs()}, L440--447); other call sites guard with \texttt{is not None} (e.g.\ L454,615,862).

\textbf{Suggested fix (diff posted in the issue).} guard the dereference, e.g. \texttt{dlg = getattr(self.parent, "\_additional\_inputs\_dialog", None)} then call \texttt{dlg.lock(...)} only if not \texttt{None}.

\subsection{\#98 --- \texttt{osdag\_core} metadata declares bogus \texttt{sqlite} dependency (T5)}
\textbf{Link:} \url{https://github.com/Nidhikhare12/OsdagBridge/issues/98}

\textbf{Statement.} The installed \texttt{osdag\_core 0.0.0} package used with OsdagBridge declares \texttt{sqlite} in its \texttt{Requires-Dist} metadata. \texttt{sqlite} is not a valid/installable PyPI distribution for the stdlib module (\texttt{sqlite3}); this can break dependency resolution in clean environments.

\textbf{Evidence.} \texttt{importlib.metadata.distribution("osdag\_core").metadata.get\_all("Requires-Dist")} includes \texttt{'sqlite'} (recorded in \texttt{evidence/logs/phase1\_installed\_meta.txt}); \texttt{tests/test\_phase1\_packaging.py::test\_osdag\_core\_declares\_bogus\_sqlite\_dep} confirms it.

\textbf{Suggested fix (posted in the issue).} remove \texttt{sqlite} from \texttt{osdag\_core} \texttt{Requires-Dist} (stdlib \texttt{sqlite3} needs no pip dependency) or reference a corrected build in OsdagBridge setup.

\subsection{\#99 --- \texttt{pyproject.toml} runtime dependencies omit required packages and declare unused \texttt{pydantic} (T5)}
\textbf{Link:} \url{https://github.com/Nidhikhare12/OsdagBridge/issues/99}

\textbf{Statement.} \texttt{pyproject.toml} declares \texttt{[scipy,fiona,shapely,pydantic,matplotlib,openseespy,opsvis,xarray]} but omits packages directly imported/used in source (e.g.\ \texttt{PySide6}, \texttt{numpy}, \texttt{osdag\_core}, \texttt{yaml} usage at \texttt{template\_page.py:812,839}, \texttt{navcube}, \texttt{OCC}, \texttt{plotly}, \texttt{ifcopenshell}, \texttt{openpyxl}, \texttt{pandas}), while \texttt{pydantic} has no observed imports in \texttt{src}. With \texttt{README} instructing \texttt{pip install -e .}, affected features can fail in fresh installs.

\textbf{Evidence.} Import/usage proof in \texttt{evidence/logs/phase1\_imports\_clean.txt} and \texttt{evidence/logs/phase1\_p1\_002\_proof.txt}; \texttt{tests/test\_phase1\_packaging.py::test\_pyproject\_missing\_key\_runtime\_deps}.

\textbf{Suggested fix (posted in the issue).} add missing runtime deps and/or optional extras to match actual requirements, reconcile \texttt{pydantic}, and align manifests/README.

\subsection{\#100 --- inconsistent dependency declarations across manifests (T6)}
\textbf{Link:} \url{https://github.com/Nidhikhare12/OsdagBridge/issues/100}

\textbf{Statement.} \texttt{pyproject.toml}, \texttt{requirements.txt}, and \texttt{environment.yml} declare different package sets (e.g.\ \texttt{PySide6/numpy/ifcopenshell/openpyxl} present in requirements but not in pyproject deps; \texttt{matplotlib/openseespy/opsvis} differ; env relies on \texttt{osdag::osdag} and \texttt{-e .}), creating conflicting install guidance.

\textbf{Evidence.} Manifests read from \texttt{dev @ d74f56d}; compared against actual imports/usage (\texttt{evidence/logs/phase1\_*}).

\textbf{Suggested fix (posted in the issue).} establish a single source of truth (e.g.\ \texttt{pyproject} + extras), align/generate other manifests, and document per-workflow install paths.

\subsection{\#101 --- README documents CLI usage but CLI entrypoint is a stub that ignores arguments (T6)}
\textbf{Link:} \url{https://github.com/Nidhikhare12/OsdagBridge/issues/101}

\textbf{Statement.} \texttt{README.md} presents CLI support and documents example commands (e.g.\ \texttt{osdagbridge analyze project.yaml --solver native}, \texttt{osdagbridge report project.yaml report.pdf}), but the installed console script \texttt{osdagbridge-cli} (\texttt{pyproject.toml:85}) resolves to \texttt{osdagbridge.cli.\_\_main\_\_.main} which only prints \texttt{"osdagbridge CLI (stub)"} and ignores all arguments including \texttt{--help}.

\textbf{Evidence.} \texttt{README.md:4,18,102,108,154}; \texttt{src/osdagbridge/cli/\_\_main\_\_.py} stub; invocations recorded in \texttt{evidence/logs/phase1\_entrypoints.txt}; \texttt{tests/test\_phase1\_packaging.py::test\_cli\_stub\_ignores\_args}.

\textbf{Suggested fix (posted in the issue).} implement argument parsing and documented commands, or clearly mark CLI as unimplemented/WIP in README and align documented vs installed command names.

\subsection{\#127 --- verification example YAMLs are empty placeholders that fail validation (T6)}
\textbf{Link:} \url{https://github.com/Nidhikhare12/OsdagBridge/issues/127}

\textbf{Statement.} \texttt{examples/plate\_girder/verification\_case\_01.yaml} and \texttt{verification\_case\_02.yaml} are empty/comment-only placeholders (no valid input dictionary) and fail \texttt{validate\_osi\_inputs}.

\textbf{Evidence.} \texttt{evidence/logs/phase1\_examples.txt}: both load as \texttt{None}/no keys with \texttt{is\_valid=False}, error \texttt{['File does not contain a valid input dictionary.']}; \texttt{tests/test\_phase1\_packaging.py::test\_examples\_parse\_and\_validate}.

\textbf{Suggested fix (posted in the issue).} populate with valid representative inputs that pass validation, or clearly mark/exclude as placeholders.

\subsection{\#128 --- \texttt{example\_basic.yaml} uses placeholder keys but passes validation vacuously (T7)}
\textbf{Link:} \url{https://github.com/Nidhikhare12/OsdagBridge/issues/128}

\textbf{Statement.} \texttt{examples/plate\_girder/example\_basic.yaml} uses a placeholder schema (e.g.\ \texttt{project.span\_m}) with no real dotted OSI keys (e.g.\ \texttt{geometry.span}), yet \texttt{validate\_osi\_inputs} returns \texttt{is\_valid=True} with no errors/warnings because it polices only present known keys.

\textbf{Evidence.} \texttt{evidence/logs/phase1\_examples.txt}: top-level \texttt{['project']}, no real span/cw/skew dotted keys, \texttt{validate.is\_valid=True}; \texttt{tests/test\_phase1\_packaging.py::test\_examples\_parse\_and\_validate}.

\textbf{Suggested fix (posted in the issue).} provide real OSI-key example values, or clearly mark as a non-loadable placeholder/template.

\subsection{\#129 --- \texttt{validate\_material\_inputs} skips range checks for cross-bracing/end-diaphragm material keys (T2)}
\textbf{Link:} \url{https://github.com/Nidhikhare12/OsdagBridge/issues/129}

\textbf{Statement.} \texttt{BridgeInputValidator.validate\_material\_inputs} enforces per-field ranges for girder/deck material keys but has no branches for \texttt{material.cross\_bracing.*} or \texttt{material.end\_diaphragm.*}, so out-of-range CB/ED custom material values are accepted.

\textbf{Evidence.} Executed calls in \texttt{scripts/drift\_check.py} (\texttt{evidence/logs/phase3\_drift\_check.txt}): girder Fy=5000 rejected; cross-bracing/end-diaphragm Fy=5000 and sampled Fu/E/density extremes accepted (returned \texttt{None}). Code: \texttt{validator.py} $\approx$L78--172 handles only girder/deck keys.

\textbf{Suggested fix (posted in the issue).} extend per-field validation to CB/ED steel material keys (same/component-appropriate ranges) or explicitly handle/flag them.

\subsection{\#130 --- integer validators raise \texttt{OverflowError} on \texttt{inf}/\texttt{-inf} instead of a validation message (T2)}
\textbf{Link:} \url{https://github.com/Nidhikhare12/OsdagBridge/issues/130}

\textbf{Statement.} \texttt{BridgeInputValidator.validate\_additional\_inputs} raises \texttt{OverflowError: cannot convert float infinity to integer} for \texttt{inf}/\texttt{-inf} on several integer fields instead of returning a validation correction/message.

\textbf{Evidence.} Executed fuzz (\texttt{scripts/validator\_fuzz.py}, \texttt{evidence/logs/trackF\_validator\_fuzz.*}): 10 exceptions for \texttt{KEY\_TS\_NO\_OF\_GIRDERS}, \texttt{KEY\_WC\_LD\_LANE\_TABLE\_COUNT}, \texttt{KEY\_MP\_CB\_NO\_OF\_CROSS\_BRACINGS}, \texttt{KEY\_DS\_STUD\_COUNT}, \texttt{KEY\_DO\_LOAD\_CYCLES}.

\textbf{Suggested fix (posted in the issue).} reject non-finite values before integer conversion (e.g. in \texttt{\_to\_int}) and return an appropriate validation message.

\subsection{\#131 --- \texttt{saveOSI\_inputs} uses \texttt{yaml.dump} unsafe tags incompatible with \texttt{yaml.safe\_load} (T2)}
\textbf{Link:} \url{https://github.com/Nidhikhare12/OsdagBridge/issues/131}

\textbf{Statement.} \texttt{saveOSI\_inputs} writes with \texttt{yaml.dump} while \texttt{loadOSI\_inputs} reads with \texttt{yaml.safe\_load}; unsafe Python tags (e.g. numpy scalars, tuples, sets) can be saved but fail to load.

\textbf{Evidence.} Executed YAML round-trip (\texttt{evidence/logs/trackI\_yaml\_dump\_safe\_load.txt}): dump of dict containing \texttt{np.float64}, tuple, set produces \texttt{!!python/...}/\texttt{!!set}/\texttt{!!python/tuple} tags and \texttt{yaml.safe\_load} raises \texttt{ConstructorError}. Code: \texttt{template\_page.py} $\approx$793--850.

\textbf{Suggested fix (posted in the issue).} use \texttt{yaml.safe\_dump} and/or normalize types (numpy scalars, tuples, sets) before saving for safe round-trip.

\subsection{\#132 --- skew sign is not mirror-symmetric (T2)}
\textbf{Link:} \url{https://github.com/Nidhikhare12/OsdagBridge/issues/132}

\textbf{Statement.} For a symmetric section (no median/footpath), flipping skew sign mirrors the bridge; demands should match. Executed headless designs (span 30, cw 10.5, Optimized) show +15° vs -15°: \texttt{util.shear} 59.28 vs 41.08, \texttt{util.fatigue} 183.74 vs 167.06; +7.5° vs -7.5°: 62.32 vs 42.18, 182.25 vs 175.33 — while self-weight checks are identical. Sign of a symmetric parameter changes the design.

\textbf{Evidence.} \texttt{dictionaries/OR\_SKEW\{15,M15,75,M75\}/output.json}, \texttt{evidence/logs/oracle\_runs\{1,2\}.log}; scripts \texttt{oracle\_runs.py}, \texttt{design\_runner.py}.

\textbf{Suggested fix (posted in the issue).} make live-load placement / oblique-mesh load mapping mirror-invariant; regression test that DCR magnitudes for positive and negative skew match within tolerance.

\subsection{\#133 --- \texttt{design()} non-termination on valid span-45 + median + footpath inputs (T2)}
\textbf{Link:} \url{https://github.com/Nidhikhare12/OsdagBridge/issues/133}

\textbf{Statement.} Valid in-limits inputs (span 45 m, cw 15 m, median Yes, footpaths) make \texttt{design()} exceed 600 s without producing output; a no-timeout rerun consumed $>12$ min CPU with no result. Control cases complete in 28--261 s. Non-convergent \texttt{fsolve} warnings appear in the case logs.

\textbf{Evidence.} \texttt{evidence/logs/design\_sweep\_run.log} (SW03/SW06 timeout), \texttt{evidence/logs/SW03\_rerun.log}, \texttt{dictionaries/SW03/case.log}, controls \texttt{dictionaries/OR\_SPAN45/result.json}.

\textbf{Suggested fix (posted in the issue).} bound the moving-load/envelope solver so non-convergent cases terminate with a clear warning; add a regression test with a time budget.

\subsection{\#134 --- wrong PASS: silent live-load skip gives zero ULS demand (T1)}
\textbf{Link:} \url{https://github.com/Nidhikhare12/OsdagBridge/issues/134}

\textbf{Statement.} Valid inputs (span 30, cw 7.5, skew 0, footpath Both Sides, median Yes) make governing-LL detection fail (\texttt{analyser.py:1998} warning); no live-load case is created, \texttt{mu\_applied=0.0}, all structural DCRs $=0$, yet the design reports \texttt{Status PASS}. Four near-neighbor controls (one field changed) show normal nonzero demands.

\textbf{Evidence.} \texttt{dictionaries/SW05/case.log}, \texttt{dictionaries/P5\_SW05R/run.log}, \texttt{dictionaries/P5\_SW05R/output.json}; controls \texttt{dictionaries/P5\_MED\_NOFOOT}, \texttt{P5\_FOOT\_NOMED}, \texttt{P5\_MED\_CW105}.

\textbf{Suggested fix (posted in the issue).} use \texttt{method="nearest"} on the \texttt{Element} lookup and/or raise on governing-LL detection failure instead of continuing to PASS.

\subsection{\#135 --- second \texttt{design()} on same object crashes in LTB buckling (T1)}
\textbf{Link:} \url{https://github.com/Nidhikhare12/OsdagBridge/issues/135}

\textbf{Statement.} Calling \texttt{design()} twice on the same \texttt{PlateGirderBridge} object: run 1 completes, run 2 crashes with \texttt{ValueError: math domain error} in \texttt{irc22\_2015.py:639} (\texttt{math.sqrt(beta\_b*(Zp*fy)/Mcr)}), stage 5 LTB check — stale per-run state. No duplicate (\#85 is CAD-display on 2nd run, no crash).

\textbf{Evidence.} guarded harness \texttt{dictionaries/IDEM2/same.py}; \texttt{evidence/logs/determinism\_2.log} (RUN1\_DONE then ValueError); full traceback \texttt{dictionaries/IDEM2/same.log}. An unguarded earlier harness (forkserver re-exec) was discarded as an artifact.

\textbf{Suggested fix (posted in the issue).} reset all per-run section/DTO state at the start of \texttt{design()}, guard the LTB sqrt, and add an idempotence regression test.

\subsection{\#153 --- \texttt{saveOSI\_inputs} does not enforce the .osi file extension (T3)}
\textbf{Link:} \url{https://github.com/Nidhikhare12/OsdagBridge/issues/153}

\textbf{Statement.} In \texttt{saveOSI\_inputs()} (\texttt{desktop/ui/template\_page.py:793--826}), the path returned by \texttt{QFileDialog.getSaveFileName(self, ``Save Design Inputs'', default\_dir, ``Input Files(*.osi)'', None)} is passed unmodified to \texttt{open(filePath, 'w')}. Nothing checks or appends the \texttt{.osi} extension, so a user typing \texttt{my\_design.txt} saves a file that the loader's \texttt{*.osi} filter will not list. A plain rename is required before the file can be loaded.

\textbf{Evidence.} \texttt{notes/osi\_load\_path.md} (OSI save/load path analysis, STEP 2); \texttt{template\_page.py:802--806} (dialog), \texttt{template\_page.py:812} (\texttt{open(filePath, 'w')}); we recorded the finding as \textbf{P6-015} (T3, static) on 2026-11-18 in \texttt{notes/all\_issues\_and\_bugs.md}. The identical review (same file/lines, same fix) was filed upstream as \#153 (2026-11-18); P6-015 was reclassified to T4 duplicate. Comment draft: \texttt{notes/issue\_drafts/25\_comment\_153\_osi\_extension\_t4.md}.

\textbf{Expected vs actual.} \emph{Expected:} \texttt{.osi} appended automatically or a warning shown. \emph{Actual:} any filename accepted; no enforcement, no warning.

\textbf{Suggested fix (as filed in the issue).} \texttt{if not filePath.lower().endswith(".osi"): filePath += ".osi"} before \texttt{open(...)}; or display a warning for unsupported extensions.

\textbf{Severity / confidence.} S3 (usability; the saved file is invisible to the load filter) / High (static, path unambiguous).

\subsection{\#154 --- \texttt{loadOSI\_inputs} silently swallows exceptions from \texttt{solve\_extend\_basic\_input\_dict} (T3)}
\textbf{Link:} \url{https://github.com/Nidhikhare12/OsdagBridge/issues/154}

\textbf{Statement.} In \texttt{loadOSI\_inputs()}, the order is: (1) \texttt{validate\_osi\_inputs(data)}; (2) on success, \texttt{populate\_from\_dict(data)} populates the dock; (3) \texttt{solve\_extend\_basic\_input\_dict(self.input\_dict)} runs inside a bare \texttt{try/except Exception: pass}. Any exception raised by the solve step --- invalid material values that slipped through validation, or any arithmetic failure --- is swallowed with no log and no message. The UI is left showing the loaded values even though the solver rejected them.

\textbf{Evidence.} \texttt{notes/osi\_load\_path.md} (item 4: exceptions silently swallowed); \texttt{template\_page.py:847--870} (validate $\rightarrow$ populate $\rightarrow$ \texttt{except Exception: pass}); recorded as \textbf{P6-017} (T3, static) on 2026-11-18 in \texttt{notes/all\_issues\_and\_bugs.md}. The identical review (same code path) was filed upstream as \#154 (2026-11-18); P6-017 was reclassified to T4 duplicate. Comment draft: \texttt{notes/issue\_drafts/26\_comment\_154\_osi\_swallow\_t4.md}.

\textbf{Expected vs actual.} \emph{Expected:} the solve failure surfaces an error to the user and the load aborts (or the dock is not populated). \emph{Actual:} the exception disappears; the dock shows the loaded values as if the load succeeded.

\textbf{Suggested fix (as filed in the issue).} Catch the exception explicitly and surface it (e.g. \texttt{CustomMessageBox} error), preserving diagnostics and keeping UI/internal state consistent, instead of \texttt{pass}.

\textbf{Severity / confidence.} S2 (wrong data shown as if load succeeded) / High (static, code path unambiguous).

\subsection{\#155 --- web application layer absent (stub-only \texttt{web/}; T5/T6)}
\textbf{Link:} \url{https://github.com/Nidhikhare12/OsdagBridge/issues/155}

\textbf{Statement.} The 2026-11-18 GATE 0 inventory found that the OsdagBridge repository contains no runnable web application: no Django backend, no React frontend, and no HTTP API. The repository is a desktop application (PySide; UI in \texttt{desktop/ui/template\_page.py}) whose \texttt{web/} tree at the repository root contains only stubs. Web (Django/React) testing therefore cannot start; any documentation or task that assumes a web layer misstates the repository.

\textbf{Corrected repository-root inventory (read-only).}

\begin{longtable}{@{}p{5.2cm}p{5.6cm}p{4.2cm}@{}}
\toprule
Path / file & Content & Verdict \\
\midrule
\texttt{web/backend/manage.py} & \texttt{print("Django manage stub")} & stub --- no Django project \\
\texttt{web/frontend/package.json} & \texttt{\{"name":"osdagbridge-frontend","version":"0.0.0"\}} & stub --- no scripts / dependencies \\
\texttt{web/frontend/src/App.jsx} & \texttt{export default function App()\{ return null; \}} & placeholder --- renders nothing \\
\texttt{web/frontend/src/components/.gitkeep} & \texttt{\# placeholder} & intentionally empty \\
\texttt{web/.DS\_Store} & 6,148 bytes (macOS artifact) & noise \\
\bottomrule
\end{longtable}

No \texttt{settings.py}, \texttt{urls.py}, \texttt{wsgi.py}, models or views; no \texttt{urlpatterns} anywhere in the tree; no \texttt{fetch(}/\texttt{axios(} calls; no import of \texttt{osdagbridge.core} under \texttt{web/}. Note: \#155's evidence table lists the top-level \texttt{web/} directory as ``Not found'', which is inaccurate --- the directory exists but holds stubs; the conclusion (no usable web layer) is correct.

\textbf{Evidence.} \texttt{notes/web\_inventory.md} (corrected inventory); commands \texttt{find} under the clone's \texttt{web/} and \texttt{cat} of the four stub files; \texttt{grep -rn urlpatterns} (no matches). Recorded as \textbf{P1-007} (T9) in Phase 1; reclassified to T4 duplicate of \#155 on 2026-11-18. Comment draft: \texttt{notes/issue\_drafts/27\_comment\_155\_web\_layer\_absent\_t4.md}.

\textbf{Expected vs actual.} \emph{Expected:} a runnable Django + React layer, or documentation explicitly stating the desktop-only scope. \emph{Actual:} stub files only; no runnable app; web testing requires building one.

\textbf{Suggested resolution (as filed in the issue).} If web testing is in scope: implement the Django backend, React frontend and API integration. If the project is desktop-only: update README/requirements to describe desktop architecture and remove references to a non-existent web application. Testing-scope blocker --- not a defect in the desktop application.

\textbf{Severity / confidence.} Blocker (web testing only) / High (inventory reproducible; conclusion independent of path detail).
